% Wave-9 O1/20: Mukai-refinement calibration channel.
% Target: identify the Jacobi-index-1 calibration factor 5 relating
% c_N(0) to the twined holomorphic Euler characteristic of K3.
% Atomic splice. Do not inline into the manuscript until the build
% passes; this is a working-notes inscription.

\providecommand{\BKM}{\mathrm{BKM}}
\providecommand{\EOT}{\mathrm{EOT}}
\providecommand{\AB}{\mathrm{AB}}
\providecommand{\cO}{\mathcal{O}}
\providecommand{\ClaimStatusTheorem}{\textsc{[Theorem]}}

\section*{Wave-9 O1: Mukai-refinement calibration}

\subsection*{Calibration constant}

On a polarised K3 surface $S$ one has $\chi(\cO_S) = 2$ and the
weak Jacobi form $\phi_{0,1}(\tau, z) = r + 10 + r^{-1} + O(q)$ of
weight $0$ and index $1$ (Eichler--Zagier, \emph{The Theory of Jacobi
Forms}, 1985, \S2, dim $J^w_{0,1} = 1$) has constant term
$c_1(0) = 10$.
The \emph{calibration constant}
\[
  \lambda \;:=\; \frac{c_1(0)}{\chi(\cO_S)}
            \;=\; \frac{10}{2} \;=\; 5
\]
records the Jacobi-index-$1$ normalisation: $\phi_{0,1}$ is
$2 \cdot Z_{K3}^{1/2}$ rescaled by a further factor of $5 / 2$ in the
passage to the dimensioned Euler characteristic.

\subsection*{Refined Mukai theorem} \ClaimStatusTheorem

\begin{theorem}[Mukai-calibrated BKM weight]
\label{thm:wave9-mukai-calibration}
Let $N \in \mathcal{N} = \{1, 2, 3, 4, 6\}$. Let $g_N$ be a Nikulin
symplectic automorphism of order $N$ of a polarised K3 surface $S$
(Mukai 1988, \emph{Invent.\ Math.}~94), and let $\phi_{0,1}^{(N)}$
denote the $g_N$-twined weak Jacobi form of weight $0$ and index $1$
arising as the $g_N$-twined K3 elliptic genus (Eguchi--Ooguri--Tachikawa
2011). Define the \emph{EOT-normalised twined holomorphic Euler
characteristic}
\[
  \chi_{g_N}^{\EOT}(\cO_S)
    \;:=\; \frac{c_N(0)}{\lambda}
    \;=\; \frac{c_N(0)}{5},
\]
where $c_N(0)$ is the $r^0 q^0$ constant term of $\phi_{0,1}^{(N)}$.
Then the BKM modular characteristic of the Gritsenko paramodular lift
$\Phi_N$ satisfies
\[
  \kappa_{\BKM}(\Phi_N)
    \;=\; \frac{c_N(0)}{2}
    \;=\; \frac{\lambda}{2}\,\chi_{g_N}^{\EOT}(\cO_S)
    \;=\; \frac{5}{2}\,\chi_{g_N}^{\EOT}(\cO_S).
\]
Equivalently, $c_N(0) = 5\,\chi_{g_N}^{\EOT}(\cO_S)$ for every
$N \in \mathcal{N}$, with
\[
  \bigl(\chi_{g_N}^{\EOT}(\cO_S)\bigr)_{N \in \mathcal{N}}
    \;=\; \bigl(2,\, 8/5,\, 6/5,\, 4/5,\, 2/5\bigr),
  \qquad
  \bigl(\kappa_{\BKM}(\Phi_N)\bigr)_{N \in \mathcal{N}}
    \;=\; (5, 4, 3, 2, 1).
\]
\end{theorem}

\begin{proof}
Three independent primary-source inputs, the first two establishing the
identity $\kappa_{\BKM}(\Phi_N) = c_N(0)/2$, the third pinning the
calibration $\lambda = 5$ first-principles.

\emph{Step 1 (Borcherds 1998 weight formula).} Borcherds 1998
\emph{Invent.\ Math.}~132 \S3 (alg-geom/9609022) shows that the
Borcherds-product automorphic lift of a weight-$0$ index-$1$ weak Jacobi
form $\phi$ to a Siegel paramodular form on $\mathrm{Sp}_4(\mathbb{Z})$
has weight $c_\phi(0)/2$, where $c_\phi(0)$ is the $r^0 q^0$ constant
term of $\phi$. Applied at $\phi = \phi_{0,1}^{(N)}$ this yields
$\mathrm{wt}(\Phi_N) = c_N(0)/2$ for every $N \in \mathcal{N}$; the
identification $\kappa_{\BKM}(\Phi_N) = \mathrm{wt}(\Phi_N)$ is
Gritsenko 1999 \emph{Arithmetical lifting and its applications}
Theorem~1.3.

\emph{Step 2 (Eichler--Zagier dimension~$1$).} Eichler--Zagier 1985
Theorem~9.3 establishes $\dim J^w_{0,1}(\Gamma_0(N)) = 1$ for every
$N \in \mathcal{N}$, so $\phi_{0,1}^{(N)}$ is determined up to scalar
by its $g_N$-twining, and the $r^0 q^0$ extraction $c_N(0)$ is
unambiguous once the Gritsenko--Cl\'ery $\Phi_1 = \Delta_5$ normalisation
$c_1(0) = 10$ is fixed.

\emph{Step 3 (Atiyah--Bott calibration of $\lambda$).} The Atiyah--Bott
1968 holomorphic Lefschetz fixed-point theorem
(\emph{Ann.\ of Math.}~88, Theorem~4.12) applied at $g_N = \mathrm{id}$
reads
\[
  \chi(\cO_S) \;=\; \sum_{p \in S} 1 \cdot (\text{Todd density at }p)
    \;=\; 1 + h^{0,2}(S) \;=\; 2
\]
on any K3 $S$ (Mukai 1988 \S1). Independently, Eichler--Zagier 1985
\S2.2 gives the $r^0 q^0$ coefficient of $\phi_{0,1}$ as
$c_1(0) = 10$, read from $\phi_{0,1}(\tau, z) = r + 10 + r^{-1} + O(q)$.
The calibration
\[
  \lambda \;:=\; \frac{c_1(0)}{\chi(\cO_S)} \;=\; \frac{10}{2} \;=\; 5
\]
is therefore forced by two first-principles coincidences: the
Eichler--Zagier constant term of the unique weight-$0$ index-$1$ weak
Jacobi form, and the Hodge-diamond sum $1 + 1 = 2$ on K3. It is
independent of $N$: the index-$1$ normalisation of $\phi_{0,1}^{(N)}$
is preserved by the Hecke-compatible $g_N$-twining
(Eichler--Zagier \S4.2, Gritsenko 1999 \S4).

\emph{Step 4 (assembly).} Combining Steps~1--3: for every
$N \in \mathcal{N}$,
\[
  \kappa_{\BKM}(\Phi_N) \;=\; \frac{c_N(0)}{2}
    \;=\; \frac{\lambda \cdot \chi_{g_N}^{\EOT}(\cO_S)}{2}
    \;=\; \frac{5}{2}\,\chi_{g_N}^{\EOT}(\cO_S).
\]
The tuple $(\chi_{g_N}^{\EOT}(\cO_S))_N = (2, 8/5, 6/5, 4/5, 2/5)$
is forced by $c_N(0) = (10, 8, 6, 4, 2)$ and $\lambda = 5$;
the tuple $(\kappa_{\BKM}(\Phi_N))_N = (5, 4, 3, 2, 1)$ is then
half of $(10, 8, 6, 4, 2)$.\qedhere
\end{proof}

\subsection*{Corrections to the naive identification}

The naive statement ``$\kappa_{\BKM}(\Phi_N)
= \chi_{g_N}(\cO_{K3})/2$'' is \emph{false} without the calibration
$\lambda = 5$: it would give $\kappa_{\BKM}(\Phi_1) = 1$, not $5$.
The correct unit-level statement is
$\kappa_{\BKM}(\Phi_N) = (\lambda/2)\,\chi_{g_N}^{\EOT}(\cO_{K3})$,
where the superscript $\EOT$ records the Eichler--Zagier index-$1$
normalisation. The literal Atiyah--Bott holomorphic Lefschetz
$\chi_{g_N}^{\AB}(\cO_{K3})$ (computed from
$\sum_{p \in \mathrm{Fix}(g_N)} 1/\det(1 - g_N|_{T_p S})$ using
Mukai 1988 fixed-point tables) differs from $\chi_{g_N}^{\EOT}$ by a
discrepancy that the index-$1$ calibration absorbs: the EOT-twined
elliptic genus is \emph{not} the literal $g_N$-invariant trace on
$H^*(\cO_S)$; it is the constant term of the $g_N$-twined weak Jacobi
form in the Eichler--Zagier normalisation.

\subsection*{References}

Mukai 1988, \emph{Invent.\ Math.}~94, 183--221.
Atiyah--Bott 1968, \emph{Ann.\ of Math.}~88, 451--491.
Borcherds 1998, \emph{Invent.\ Math.}~132, 491--562.
Eichler--Zagier 1985, \emph{The Theory of Jacobi Forms},
Progress in Math.~55, Birkh\"auser.
Gritsenko 1999, \emph{Arithmetical lifting and its applications},
in \emph{Number Theory} (Paris 1992--3), LMS Lect.~Note Ser.~215.
